\documentclass[11pt,a4paper]{article}
\usepackage[T1]{fontenc}
\usepackage[utf8]{inputenc}
\usepackage{lmodern}
\usepackage[english]{babel}
\usepackage[margin=18mm]{geometry}
\usepackage{booktabs,tabularx,array}
\usepackage[hidelinks]{hyperref}

\setlength{\parindent}{0pt}
\setlength{\tabcolsep}{4pt}
\renewcommand{\arraystretch}{1.15}
\newcolumntype{L}{>{\raggedright\arraybackslash}X}
\newcommand{\finding}[1]{\textbf{#1}\par\smallskip}
% Trenner zwischen Papern. Umbruch vorziehen, wenn weniger als zehn Zeilen frei sind (needspace-Trick ohne Paket).
% Linie, Abstand und Titel bleiben zusammen: \@nobreaktrue lässt \section* keinen eigenen Umbruchpunkt setzen.
% Die Linie ist Leaders-Glue, damit sie am Seitenanfang wegfällt und nur zwischen zwei Papern auf einer Seite erscheint.
\makeatletter
\newcommand{\papersep}{\par\vskip\z@\@plus 10\baselineskip\penalty-100\vskip\z@\@plus -10\baselineskip
  \vspace{4ex}\leaders\hrule height 0.4pt\vskip 0.4pt\nobreak\vspace{3.5ex}\@nobreaktrue}
\makeatother

\begin{document}

% Layout inspired by Downloads/IMG_0500.HEIC.
\begin{table}[!t]
\centering
\small
\begin{tabularx}{\textwidth}{@{}LL@{}}
\toprule
\textbf{In-kernel} & \textbf{Kernel bypass / userspace} \\
\midrule
eBPF kernel programs \cite{xdp2018,freitas2022,tu2021}\newline
\hspace*{1em}XDP (native/generic) \cite{xdp2018,tu2021}\newline
\hspace*{1em}SK\_MSG \cite{middlenet2023}\newline
Linux network stack \cite{xdp2018,tu2021,ara2020}\newline
\hspace*{1em}AF\_PACKET sockets \cite{valia2026}
&
DPDK \cite{gallenmueller2015,metronome2023,valia2026,tu2021,ara2020}\newline
Snabb \cite{ara2020}\newline
PF\_RING ZC \cite{gallenmueller2015}\newline
netmap \cite{gallenmueller2015}\newline
AF\_XDP \cite{mole2025,edge2023,valia2026,tu2021}\newline
uXDP over DPDK/AF\_XDP \cite{uxdp2025}\newline
VPP over DPDK \cite{ara2020} \\
\bottomrule
\end{tabularx}
\end{table}

\begin{table}[!t]
\centering
\footnotesize
\renewcommand{\arraystretch}{1.05}
% Spaltenbreiten: Paper schmal, Comparison basis breit, Summe der Faktoren bleibt 5.
\begin{tabularx}{\textwidth}{@{}>{\hsize=.7\hsize}L>{\hsize=1.0\hsize}L>{\hsize=1.05\hsize}L>{\hsize=.95\hsize}L>{\hsize=1.3\hsize}L@{}}
\toprule
\textbf{Paper} & \textbf{Approaches} & \textbf{Workload} &
\textbf{Metrics} & \textbf{Comparison basis} \\
\midrule
Packet IO (2015) \cite{gallenmueller2015} &
DPDK, netmap, PF\_RING ZC &
Forwarding, fixed CPU cost, random accesses to 1--512~MB working sets &
Packet rate, cycles/packet, batch/cache effects, latency at 99\% line rate &
Same CPU, 10-Gbit NICs and forwarding algorithm. \\

AF\_XDP / OpenNIC (2025) \cite{mole2025} &
Mixed-mode AF\_XDP zero-copy, DPDK on OpenNIC &
Packet drop and L2 MAC-swap forwarding &
Rate vs.\ packet size and queue count &
Same FPGA, different drivers and benchmark tools. \\

XDP (2018) \cite{xdp2018} &
Native XDP, DPDK \texttt{testpmd}, Linux stack &
64-byte drop and MAC-swap forwarding, core scaling &
Rate vs.\ cores, CPU use vs.\ load, forwarding latency &
Same testbed. RSS selects cores. Linux stack only in drop/CPU tests. \\

Linux survey (2022) \cite{freitas2022} &
Hardware, software, virtualisation &
Offload, RDMA, GPU, zero-copy, bypass, XDP, virtual NICs &
Resource efficiency, packet rate, security, flexibility &
Literature review, no shared benchmark/testbed. \\

uXDP (2025) \cite{uxdp2025} &
Native/generic XDP, userspace eBPF via DPDK/AF\_XDP &
Counters, filtering, load balancing, Katran, MAC-swap echo (\texttt{xdp\_tx}) &
NF throughput, unloaded RTT, compiler/inlining effects &
Same programs/hardware, one core. Mode comparison: optimised source LLVM IR. \\

Metronome (2023) \cite{metronome2023} &
DPDK busy polling, adaptive sleep/wake, native XDP &
L3 \texttt{l3fwd}, load changes, concurrent CPU workloads &
Rate, loss, CPU use, latency vs.\ sleep interval, power &
DPDK/Metronome: same \texttt{l3fwd}. Different thread/core allocations. \\

MiddleNet (2023) \cite{middlenet2023} &
DPDK (PMD, RTE\_RING), AF\_XDP, eBPF SK\_MSG, SR-IOV &
L3 routing, L2 forwarding, HTTP proxy, URL rewrite, chain scaling &
MLFR, CPU and latency at MLFR, HTTP request rate/latency, copy and interrupt overhead &
Same workloads/testbed. Five 60-s L2/L3 runs. Data paths and CPU budgets differ. \\

Edge data centers (2023) \cite{edge2023} &
Native XDP, AF\_XDP (busy loop, syscall, \texttt{poll()}), XDP with AF\_XDP sockets &
Drop, transit, local traffic, CPU/memory scaling, firewall, load balancer, memcached &
Throughput, CPU by execution context, LLC hits/misses &
40-Gbit testbed, usually one core. 10 runs (mean/SD). Drop/transit loss below 0.1\%. \\

UPF I/O modes (2026) \cite{valia2026} &
AF\_PACKET, AF\_XDP zero-copy, CNDP, DPDK PMD as port drivers of one BESS-UPF &
5G UPF pipeline (GTP-U, PDR/FAR/QER), 64--1518~B, 1--8 workers, 24-h soak &
Non-drop rate (RFC~2544), in-pipeline latency percentiles, per-core rate, worker scaling &
Same host, NIC, kernel and UPF image, only the port driver changes. Different generators per mode. AF\_XDP/CNDP capped at two workers. \\

OVS dataplane (2021) \cite{tu2021} &
OVS kernel module, OVS-DPDK, OVS-AF\_XDP, eBPF/XDP variants &
NSX rule set with TCP, 64-byte P2P/PVP/PCP forwarding, \texttt{netperf} RR, 1--6 queues &
TCP Gbps, lossless Mpps, CPU by context, P50/P90/P99 latency &
Same switch and testbed, I/O backend swapped. Kernel path uses RSS over all cores, userspace paths pinned PMD threads. \\

Inter-container bypass (2020) \cite{ara2020} &
Linux bridge/sockets, DPDK basicfwd, OVS-DPDK, VPP, Snabb, SR-IOV &
Container pairs on one or two hosts, 64--1500~B, bursts 4--256, 1--8 pairs &
Max receive rate, total throughput vs.\ pairs, mean round-trip latency &
Same hosts and benchmark apps (\texttt{nfv-testperf}), polling, fixed CPU clock. Switching logic differs per switch. \\
\bottomrule
\end{tabularx}
\end{table}

\null % Keep both tables at the top instead of centring them on a float-only page.
\clearpage

\begingroup
\small
\section*{\large\normalfont Comparison of Frameworks for High-Performance Packet IO~\cite{gallenmueller2015}}
\finding{DPDK and PF\_RING ZC outperform netmap, and the CPU cycles available per packet bound throughput.}
DPDK and PF\_RING ZC process packets in user space using zero-copy.
The Linux stack handles packets in kernel space with security and robustness checks.
netmap combines userspace processing with kernel buffer checks and packet transfer through system calls.
DPDK and PF\_RING ZC outperform netmap in most tested forwarding scenarios.
Available CPU cycles limit throughput.

\subsection*{\normalsize\normalfont Metrics}
\begin{itemize}
\setlength{\itemsep}{2pt}
\item \emph{CPU cycles per packet.}
More cycles per packet reduce throughput (Figure~2a).
$C_{\mathrm{task}}$ is the application's processing cost per packet.

\item \emph{Cache size vs.\ throughput.}
Throughput drops and cache misses rise when the working set reaches the CPU's 4~MB L3 cache (Figures~3a and~3b).

\item \emph{Batch size, throughput and latency.}
Packets per NIC batch affect throughput and CPU cycles.
Larger batches increase queueing delay.
Very small batches sharply reduce throughput.
netmap needs large batches for high throughput.
Latency is measured at 99\% of line rate.
\end{itemize}

\subsection*{\normalsize\normalfont Workload}
\begin{itemize}
\setlength{\itemsep}{2pt}
\item \emph{Simulated task $C_{\mathrm{task}}$.}
Fixed CPU cycles per packet simulate application work.
\item \emph{Memory access.}
One random memory access per packet, with 1--512~MB working sets.
\end{itemize}

\par
\endgroup

\papersep
\begingroup
\small
\section*{\large\normalfont Performance Implications at the Intersection of AF\_XDP and Programmable NICs~\cite{mole2025}}
\finding{DPDK beats mixed-mode AF\_XDP only for 64-byte packets, while both perform nearly equally for larger packets.}
The OpenNIC driver adds eBPF/XDP loading and AF\_XDP zero-copy.
Without driver support, AF\_XDP uses copy mode.
Focus: splitting AF\_XDP functionality across the FPGA, driver and kernel.

\subsection*{\normalsize\normalfont Processing paths and motivation}
\begin{itemize}
\setlength{\itemsep}{2pt}
\item \emph{Native XDP.}
XDP runs in the kernel's NIC driver.
Packets cross PCIe into host memory before XDP processes them.
High packet rates can exhaust host CPU capacity.
\item \emph{NIC offloading.}
NIC-side filtering or processing can reduce host CPU work.
Programmable hardware and hardware development are costly.
AF\_XDP hardware features consume FPGA resources otherwise available to application logic.
\item \emph{Studied hardware/kernel split.}
Kernel eBPF selects the receiving AF\_XDP socket.
Software redirection reduces hardware flow-steering requirements but adds host processing.
Further parameters: TX NAPI budget and NAPI placement on CPU cores.
\end{itemize}

AF\_XDP mixed mode shares the NIC between bypass applications and traditional Linux sockets.
Exclusive DPDK deployments assign the NIC to one application.

\subsection*{\normalsize\normalfont Metrics}
\begin{itemize}
\setlength{\itemsep}{2pt}
\item \emph{Throughput vs.\ packet size.}
Small packets require higher packet rates at the same data rate.
DPDK outperforms mixed-mode AF\_XDP (XSK) for 64-byte packets.
AF\_XDP and DPDK perform nearly equally for larger packets (Figures~7 and~8).
\end{itemize}

\subsection*{\normalsize\normalfont Workload}
\begin{itemize}
\setlength{\itemsep}{2pt}
\item \emph{Packet drop.}
Receive and drop packets.
\item \emph{L2 swap and forward.}
Swap source/destination MAC addresses and forward.
\end{itemize}

\par
\endgroup

\papersep
\begingroup
\small
\section*{\large\normalfont The eXpress Data Path: Fast Programmable Packet Processing in the Operating System Kernel~\cite{xdp2018}}
\finding{XDP drops about five times more packets per core than the Linux stack with load-proportional CPU use, while DPDK stays faster at 100\,\% CPU.}
XDP executes eBPF in the receive path before the Linux network stack.
General-purpose stacks prioritise flexibility over per-packet cost.
XDP reduces this overhead while preserving application isolation and kernel security.
DPDK remains faster in some tested scenarios.
\begin{itemize}
\setlength{\itemsep}{2pt}
\item Cooperates with the regular network stack.
\item Remains transparent to host applications.
\item CPU usage scales with load without exclusively reserved cores.
\end{itemize}

\subsection*{\normalsize\normalfont Metrics}
\begin{itemize}
\setlength{\itemsep}{2pt}
\item \emph{Packet drop rate vs.\ number of cores.}
Minimal drop processing measures system overhead and gives an upper bound on application throughput (Figure~3).
DPDK and XDP scale approximately linearly with the number of cores until they approach the PCI bus limit.
Single-core XDP achieves about five times the drop rate of the Linux stack in raw mode.

\item \emph{CPU usage vs.\ offered packet load.}
Single-core drop test (Figure~4): DPDK busy polling uses 100\% CPU.
XDP and Linux CPU usage scale with offered load.

\item \emph{Packet forwarding throughput and latency.}
Throughput vs.\ core count (Figure~5).
Latency at low and high packet rates (Table~1).
\end{itemize}

\subsection*{\normalsize\normalfont Workload}
\begin{itemize}
\setlength{\itemsep}{2pt}
\item \emph{Minimum-sized packets.}
64-byte packets maximise packet rate at a given line rate.
\item \emph{Packet drop.}
Receive and drop without application processing.
\item \emph{L2 swap and forward.}
Swap source/destination MAC addresses and forward.
Minimal rewriting gives an upper bound on forwarding throughput.
\end{itemize}

\subsection*{\normalsize\normalfont Comparison basis}
Same testbed and basic packet tasks.
DPDK uses \texttt{testpmd}.
Hardware Receive Side Scaling (RSS) steers XDP/Linux traffic to the selected cores.
The Linux stack appears in drop and CPU tests.
Minimal forwarding excludes it because it requires bridging or routing lookups.
\par
\endgroup

\papersep
\begingroup
\small
\section*{\large\normalfont A survey on accelerating technologies for fast network packet processing in Linux environments~\cite{freitas2022}}
\finding{Survey that classifies accelerators by resource use, packet rate, security and flexibility, without a common benchmark or performance ranking.}
Taxonomy of Linux packet-processing bottlenecks and accelerators: hardware, software and virtualisation (Figure~2).

\subsection*{\normalsize\normalfont Metrics}
\begin{itemize}
\setlength{\itemsep}{2pt}
\item \emph{Host resource usage and packet rate.}
Published CPU usage and packet rates.
Continuous polling trades CPU occupation at low load for higher throughput.
\item \emph{Security and flexibility.}
Kernel isolation, programmability, hardware dependencies and expandability.
\end{itemize}

\subsection*{\normalsize\normalfont Workload}
\begin{itemize}
\setlength{\itemsep}{2pt}
\item \emph{Hardware acceleration.}
TCP offload engines, programmable NICs, RDMA and GPU packet processing.
\item \emph{Software acceleration.}
Zero-copy, kernel bypass, in-kernel processing and userspace TCP.
Examples: netmap, DPDK, PF\_RING ZC and XDP/eBPF (in-kernel).
\item \emph{Virtualised networking.}
NIC passthrough, SR-IOV, VirtIO/VHost and userspace processing (VPP).
\end{itemize}

\subsection*{\normalsize\normalfont Comparison basis}
Literature synthesis without a common benchmark or testbed.
The survey's Table~2 maps studies to the four criteria.
Supports technology selection, with no overall performance ranking.
\par
\endgroup

\papersep
\begingroup
\small
\section*{\large\normalfont uXDP: Frictionless XDP Deployments in Userspace~\cite{uxdp2025}}
\finding{DPDK-mode uXDP is fastest, up to $3.3\times$ in-kernel XDP, but the figures mix I/O-path and compiler effects.}
uXDP runs unchanged XDP/eBPF programs in userspace via DPDK or AF\_XDP.
It extends \texttt{bpftime}, retaining eBPF verification and existing control-plane compatibility.
Optimisations: helper/map inlining and SIMD instructions.
DPDK mode avoids interrupt-based retrieval and enables execution where kernel eBPF is unavailable.

\subsection*{\normalsize\normalfont Deployment modes}
\begin{itemize}
\setlength{\itemsep}{2pt}
\item \emph{DPDK mode.}
Userspace RX/TX rings and packet processing bypass the kernel network stack.
Kernel-bound traffic requires reinjection.
\item \emph{AF\_XDP mode.}
Kernel XDP redirects packets to an AF\_XDP socket for userspace eBPF processing.
Shared RX/TX rings and packet buffers avoid per-packet copies in zero-copy mode.
\end{itemize}

In AF\_XDP mode, kernel XDP can pass selected packets directly to the Linux stack before userspace processing.
After uXDP processing, the evaluated version supports only \texttt{XDP\_DROP} and \texttt{XDP\_TX}.
Kernel reinjection for \texttt{XDP\_PASS} is future work (Section~7).

\subsection*{\normalsize\normalfont Metrics}
\begin{itemize}
\setlength{\itemsep}{2pt}
\item \emph{Throughput across deployment modes.}
Modes in Figure~4: DPDK (\texttt{dpdk\_llvm}), AF\_XDP (\texttt{afxdp\_llvm}), native XDP (\texttt{drv\_mode}), generic XDP (\texttt{skb\_mode}).
Up to $3.3\times$ throughput over in-kernel execution and up to 40\% higher throughput for Meta's Katran.
DPDK mode is fastest across tested functions.
AF\_XDP mode can exceed native XDP for complex functions.
Generic XDP is slowest.

\item \emph{Unloaded round-trip latency.}
\texttt{xdp\_tx} MAC-swap RTT (Figure~5): DPDK mode has the lowest latency, AF\_XDP mode slightly higher latency than native XDP.
\end{itemize}

\subsection*{\normalsize\normalfont Workload}
\begin{itemize}
\setlength{\itemsep}{2pt}
\item \emph{XDP network functions.}
Packet counters, packet-length statistics, filtering, simple load balancing and Katran, with different helper/map combinations.
\item \emph{MAC swap and transmit.}
\texttt{xdp\_tx}: swap source/destination MAC addresses and transmit back (latency benchmark).
\end{itemize}

\subsection*{\normalsize\normalfont Comparison basis}
Same programs and testbed, one CPU core.
64-byte TCP or 128-byte ICMP packets.
Existing eBPF binaries run without recompilation.
Figure~4 uses optimised source LLVM IR: results combine deployment-mode and compiler effects.
\par
\endgroup
\newpage
\papersep
\begingroup
\small
\section*{\large\normalfont Metronome: Adaptive and Precise Intermittent Packet Retrieval in DPDK~\cite{metronome2023}}
\finding{Adaptive sleep/wake keeps DPDK's throughput and loss while trading CPU time and power against buffering latency.}
DPDK busy polling occupies reserved cores at 100\% even at low traffic, consuming CPU time and power.
Metronome adapts sleep/wake intervals to packet load.
Alternating threads release CPU time for other applications.
Like XDP, it targets load-proportional CPU utilisation.

\subsection*{\normalsize\normalfont Metrics}
\begin{itemize}
\setlength{\itemsep}{2pt}
\item \emph{Latency and CPU usage vs.\ sleep interval.}
Target sleep interval varied at 5 and 10~Gbps (Figure~5).
Longer sleep intervals reduce CPU usage but increase packet buffering delay.
Adaptive tuning targets a configurable latency as load changes.

\item \emph{Latency vs.\ number of threads.}
More threads sharing a receive queue can increase latency and variance through coordination and longer sleep periods (Figures~7 and~8).

\item \emph{Power consumption.}
Sleeping can reduce polling power costs.
Savings depend on traffic load and CPU power governor.

\item \emph{Throughput and packet loss.}
Tuned Metronome achieves comparable throughput and no substantial packet-loss difference from standard DPDK.
\end{itemize}

\subsection*{\normalsize\normalfont Workload}
\begin{itemize}
\setlength{\itemsep}{2pt}
\item \emph{L3 forwarding.}
DPDK \texttt{l3fwd} with changing traffic rates.
\item \emph{CPU sharing.}
CPU-intensive applications run alongside packet I/O to use released CPU time.
\end{itemize}

\subsection*{\normalsize\normalfont Comparison basis}
DPDK/Metronome: same forwarding application and hardware.
Metronome typically uses three threads, DPDK one polling thread per receive queue.
Native XDP comparison: \texttt{xdp\_router\_ipv4} (Section~V-D).
Core counts and frame sizes differ in that comparison.
\par
\endgroup

\papersep
\begingroup
\small
\section*{\large\normalfont MiddleNet: A Unified, High-Performance NFV and Middlebox Framework With eBPF and DPDK~\cite{middlenet2023}}
\finding{DPDK paths win for small packets and event-driven eBPF paths use less CPU for large ones, but the comparison covers whole architectures rather than isolated I/O paths.}
MiddleNet combines L2/L3 network functions and L4/L7 middleboxes on one host.
L2/L3 uses DPDK kernel bypass for zero-copy delivery and high packet rates.
L4/L7 combines event-driven eBPF with the kernel protocol stack.
SR-IOV allows both paths to coexist.

\smallskip
Separate NFV/middlebox platforms can place functions on different nodes.
Service function chaining (SFC) transfers can add interrupts, copies, context switches, protocol processing and serialization/deserialization.
Overhead grows with chain length and can reduce throughput.
MiddleNet stores packet data in shared memory and passes descriptors between functions on the same host, avoiding repeated payload copies.

\subsection*{\normalsize\normalfont Shared-memory communication}
NIC access: DPDK PMD or AF\_XDP.
Descriptor transfer between functions: RTE\_RING or SK\_MSG.
L4/L7 uses sockets and a userspace message broker with the kernel protocol stack.
\begin{itemize}
\setlength{\itemsep}{2pt}
\item \emph{Event-driven eBPF.}
Negligible overhead without events.
Programs have size limits and must run to completion.
Interrupts and context switches limit performance at high packet rates.
\item \emph{Polling-based DPDK.}
PMD and RTE\_RING avoid interrupts for retrieval and descriptor transfer.
Continuous polling occupies CPU cores.
This path provides zero-copy, interrupt-free chaining.
\end{itemize}

\subsection*{\normalsize\normalfont Metrics}
\begin{itemize}
\setlength{\itemsep}{2pt}
\item \emph{Packet rate, CPU usage and latency.}
Maximum loss-free rate (MLFR), CPU utilisation and end-to-end latency at MLFR vs.\ packet size (Figure~7).
DPDK performs better for small packets.
Event-driven variants use less CPU for larger packets.
\item \emph{Application performance.}
HTTP requests/s, response latency and CPU usage vs.\ client concurrency (Figure~10).
\item \emph{Communication overhead.}
Data copies, interrupts and context switches along the chain.
\end{itemize}

\subsection*{\normalsize\normalfont Workload}
\begin{itemize}
\setlength{\itemsep}{2pt}
\item \emph{L2/L3 chain.}
L3 routing followed by L2 forwarding updates IP/MAC addresses.
DPDK: PMD and RTE\_RING.
eBPF: SK\_MSG with AF\_XDP polling or interrupts.
\item \emph{L4/L7 chain.}
HTTP reverse-proxy load balancing followed by URL rewriting.
Chain length and compute demand varied separately.
\end{itemize}

\subsection*{\normalsize\normalfont Comparison basis}
Same testbed and comparable chains, with OVS and kernel-bridge baselines.
L2/L3: mean of five 60-second runs.
Results compare complete architectures with different data paths, communication mechanisms and CPU budgets.
\par
\endgroup

\papersep
\begingroup
\small
\section*{\large\normalfont Comparing User Space and In-Kernel Packet Processing for Edge Data Centers~\cite{edge2023}}
\finding{Native XDP wins for cache-resident dropping and syscall-driven AF\_XDP for forwarding and memory-heavy processing, with latency left unmeasured.}
Native XDP vs.\ zero-copy AF\_XDP on edge servers sharing resources with applications.
Kernel-provided XDP and AF\_XDP, no custom modules.
AF\_XDP modes: busy loop, syscall-triggered driver polling and \texttt{poll()}.
Additional modes: XDP processing with AF\_XDP sockets enabled.

\subsection*{\normalsize\normalfont Metrics}
\begin{itemize}
\setlength{\itemsep}{2pt}
\item \emph{Throughput.}
RFC2544-based input-rate adjustment until drop/transit loss is below 0.1\%.
XDP performs well for cache-resident dropping.
Syscall-driven AF\_XDP benefits forwarding and memory-intensive processing.
\item \emph{CPU usage.}
Userspace, system-call and software-interrupt CPU usage.
\item \emph{Cache and memory behaviour.}
LLC load/store hits and misses vs.\ working set and processing cost.
Latency evaluation is left to future work.
\end{itemize}

\subsection*{\normalsize\normalfont Workload}
\begin{itemize}
\setlength{\itemsep}{2pt}
\item \emph{Traffic classes.}
Drop, pass-through to remote destinations, local delivery through the kernel stack.
\item \emph{Processing and memory demand.}
Minimal MAC swap isolates packet I/O overhead.
Repeated checksums add CPU work.
Random accesses to growing arrays add memory demand.
\item \emph{Application examples.}
Firewall, load balancer, local memcached and mixed traffic.
\end{itemize}

\subsection*{\normalsize\normalfont Comparison basis}
Same 40-Gbps testbed, usually one CPU core.
Ten repetitions, mean and standard deviation.
Results depend on traffic class, memory demand and driver scheduling mode.
\par
\endgroup

\papersep
\begingroup
\small
\section*{\large\normalfont Comparing UPF Dataplane I/O Modes in a Cloud-Native 5G Core: AF\_PACKET, AF\_XDP, CNDP, and DPDK on SD-Core BESS-UPF~\cite{valia2026}}
\finding{At two workers AF\_XDP, CNDP and DPDK are within measurement noise and AF\_XDP leads per core, DPDK's lead is a scaling-ceiling effect, and AF\_PACKET is an order of magnitude slower.}
One open-source 5G user plane function (SD-Core BESS-UPF) runs under four I/O modes.
Only the BESS port driver changes: AF\_PACKET and AF\_XDP through DPDK poll-mode drivers, CNDP through its own port type, DPDK through a native PMD over \texttt{vfio-pci}.
Deployment requirements, operational pitfalls and a decision guide are reported next to the measurements (Sections~3, 7 and~9).

\subsection*{\normalsize\normalfont Metrics}
\begin{itemize}
\setlength{\itemsep}{2pt}
\item \emph{Throughput.}
Non-drop rate per RFC~2544 over 64--1518~B (Figure~18, Tables~3 and~4).
At 64~B: CNDP 5.52 and AF\_XDP 6.47~Mpps with two workers, DPDK 10.30~Mpps with four and 13.09~Mpps with eight; AF\_PACKET 0.25~Mpps at 86~B.
\item \emph{Latency and jitter.}
In-pipeline probes give mean, P50, P99 and P99.9 (Table~5): DPDK 8.1~\textmu s mean, CNDP and AF\_XDP 13--14~\textmu s.
A second AF\_XDP worker adds P99 jitter through contention on the kernel redirect path (Table~6).
\item \emph{CPU and core scaling.}
Per-core rate and scaling efficiency (Table~7): AF\_XDP 3.69, DPDK 2.99, CNDP 2.65~Mpps per core.
A per-netdev AF\_XDP socket limit pins CNDP and AF\_XDP at two workers.
\end{itemize}

\subsection*{\normalsize\normalfont Workload}
\begin{itemize}
\setlength{\itemsep}{2pt}
\item \emph{UPF pipeline.}
GTP-U decapsulation, PDR lookup, QER metering and FAR action with a single rule set.
\item \emph{Synthetic and real traffic.}
\texttt{testpmd} and kernel \texttt{pktgen} for I/O ceilings; a 24-hour soak and a real-RAN run validate CNDP only.
\end{itemize}

\subsection*{\normalsize\normalfont Comparison basis}
Same host, NIC (Intel XXV710), real-time kernel, Kubernetes cluster and UPF image; the port driver is the only variable (Table~2).
Caveats: a different load generator per mode, traffic hairpinned inside the NIC because of 1-GbE optics, a NIC receive path that caps every mode near 3~Mpps per core, no energy measurement.
\par
\endgroup

\papersep
\begingroup
\small
\section*{\large\normalfont Revisiting the Open vSwitch Dataplane Ten Years Later~\cite{tu2021}}
\finding{DPDK stays ahead of AF\_XDP on physical and VM paths at lower CPU cost, the kernel datapath keeps up only by spending about eight cores through RSS, and AF\_XDP wins container forwarding via XDP redirect.}
Open vSwitch (OVS) moves from a kernel-module datapath to userspace over AF\_XDP, with DPDK as the performance reference.
The motivation is maintainability (Section~2), the paper closes with two years of lessons learned (Section~6).
An eBPF datapath at the \texttt{tc} hook was rejected as 10--20\,\% slower than the kernel module (Figure~2).

\subsection*{\normalsize\normalfont Metrics}
\begin{itemize}
\setlength{\itemsep}{2pt}
\item \emph{Forwarding rate and CPU use.}
Maximum lossless rate for 64-byte packets with 1 and 1,000 flows; CPU in hyperthread units split into system, softirq, guest and user (Figure~9, Table~4).
With vhostuser, AF\_XDP is about 30\,\% slower than DPDK.
\item \emph{Latency.}
\texttt{netperf} TCP\_RR P50/P90/P99 (Figures~10 and~11): inter-host DPDK 36/38/45~\textmu s, AF\_XDP 39/41/53, kernel 58/68/94; between containers DPDK 81/136/241 versus about 15/16/20 for kernel and AF\_XDP.
\item \emph{Scaling and XDP cost.}
Throughput over 1--6 queues at 25~Gbps (Figure~12).
Single-core XDP rate falls from 14 to 4.7~Mpps as program complexity grows (Table~5).
\end{itemize}

\subsection*{\normalsize\normalfont Workload}
\begin{itemize}
\setlength{\itemsep}{2pt}
\item \emph{Production rule set.}
About 103,000 NSX OpenFlow rules with Geneve tunnels and connection tracking, TCP via \texttt{iperf} for VM and container scenarios (Figure~8).
\item \emph{Raw forwarding.}
TRex 64-byte traffic through physical-to-physical, physical-VM-physical and physical-container-physical loops.
\end{itemize}

\subsection*{\normalsize\normalfont Comparison basis}
OVS 2.14 with swapped I/O backend on one testbed.
The kernel datapath spreads load over all 16 hyperthreads via RSS, userspace datapaths use pinned PMD threads.
Virtual device type (tap versus vhostuser) and the missing AF\_XDP checksum offload (estimated) shift results.
DPDK is absent from the NSX scenario.
\par
\endgroup

\papersep
\begingroup
\small
\section*{\large\normalfont Comparative Evaluation of Kernel Bypass Mechanisms for High-performance Inter-container Communications~\cite{ara2020}}
\finding{Every bypass path beats kernel sockets by over four times, SR-IOV leads throughput and large-burst latency, VPP nearly matches raw DPDK forwarding, and Snabb trails on both throughput and latency.}
An open-source framework (\texttt{nfv-testperf}) benchmarks virtual switches between LXC containers.
Compared: Linux bridge with POSIX sockets, the DPDK basic forwarding sample, OVS-DPDK, VPP, Snabb as a learning switch, and SR-IOV with the NIC's embedded switch.
Benchmark applications poll and use one API over POSIX or DPDK.

\subsection*{\normalsize\normalfont Metrics}
\begin{itemize}
\setlength{\itemsep}{2pt}
\item \emph{Throughput.}
Maximum receive rate in Mpps and Gbps versus packet size (Figure~4).
Kernel sockets reach 0.34--0.44~Mpps, DPDK-based switches at least 2~Mpps (Table~2).
Burst sizes 32--256 did not change throughput.
\item \emph{Scalability.}
Total throughput for 1--8 container pairs on one or two hosts (Figures~4e--4h).
\item \emph{Latency.}
Mean round-trip time at low rate versus packet and burst size (Figure~5).
Only virtio-based switches reach single-digit microseconds locally; Snabb stays above 60~\textmu s.
\end{itemize}

\subsection*{\normalsize\normalfont Workload}
\begin{itemize}
\setlength{\itemsep}{2pt}
\item \emph{Sender/receiver pairs.}
Unidirectional synthetic traffic, 64--1500-byte packets, 1--20~Mpps offered, one-minute runs.
\item \emph{Client/server pairs.}
Echoed bursts of 4--128 packets for round-trip latency.
\end{itemize}

\subsection*{\normalsize\normalfont Comparison basis}
Two identical hosts (Xeon E5-2640 v4, Intel X710 10~GbE), frequency scaling and hyperthreading off; DPDK 19.05, OVS 2.11, Snabb 2019.01, VPP 19.08.
Switching logic differs: the DPDK sample does no switching, Snabb's L2 component is its bottleneck, SR-IOV switches in hardware.
Snabb runs single-host only; CPU cores per switch are not varied.
\par
\endgroup

\clearpage
\begin{thebibliography}{11}
\footnotesize
\setlength{\itemsep}{3pt}

\bibitem{gallenmueller2015}
S.~Gallenmüller et al.
\emph{Comparison of Frameworks for High-Performance Packet IO}.
ANCS, 2015.
\href{https://www.net.in.tum.de/fileadmin/bibtex/publications/papers/gallenmueller_ancs2015.pdf}{PDF}.

\bibitem{mole2025}
M.~Molè et al.
\emph{Performance Implications at the Intersection of AF\_XDP and Programmable NICs}.
eBPF Workshop, 2025.
\href{https://doi.org/10.1145/3748355.3748359}{DOI}.

\bibitem{xdp2018}
T.~Høiland-Jørgensen et al.
\emph{The eXpress Data Path: Fast Programmable Packet Processing in the Operating System Kernel}.
CoNEXT, 2018.
\href{https://doi.org/10.1145/3281411.3281443}{DOI}.

\bibitem{freitas2022}
E.~Freitas et al.
\emph{A survey on accelerating technologies for fast network packet processing in Linux environments}.
Computer Communications, 2022.
\href{https://doi.org/10.1016/j.comcom.2022.10.003}{DOI}.

\bibitem{uxdp2025}
Y.~Zheng, P.~Gavriil and M.~Kogias.
\emph{uXDP: Frictionless XDP Deployments in Userspace}.
eBPF Workshop, 2025.
\href{https://doi.org/10.1145/3748355.3748360}{DOI}.

\bibitem{metronome2023}
M.~Faltelli et al.
\emph{Metronome: Adaptive and Precise Intermittent Packet Retrieval in DPDK}.
IEEE/ACM Transactions on Networking, 2023 (journal edition).
\href{https://doi.org/10.1109/TNET.2022.3208799}{DOI}.

\bibitem{middlenet2023}
S.~Qi et al.
\emph{MiddleNet: A Unified, High-Performance NFV and Middlebox Framework With eBPF and DPDK}.
IEEE Transactions on Network and Service Management, 2023.
\href{https://doi.org/10.1109/TNSM.2023.3256891}{DOI}.

\bibitem{edge2023}
F.~Parola, R.~Procopio, R.~Querio and F.~Risso.
\emph{Comparing User Space and In-Kernel Packet Processing for Edge Data Centers}.
ACM SIGCOMM Computer Communication Review, 2023.
\href{https://doi.org/10.1145/3594255.3594257}{DOI}.

\bibitem{valia2026}
S.~Valia, N.~Rajput, A.~Gairola and V.~Rathi.
\emph{Comparing UPF Dataplane I/O Modes in a Cloud-Native 5G Core: AF\_PACKET, AF\_XDP, CNDP, and DPDK on SD-Core BESS-UPF}.
arXiv:2608.14589, 2026.
\href{https://arxiv.org/abs/2608.14589}{arXiv}.

\bibitem{tu2021}
W.~Tu, Y.-H.~Wei, G.~Antichi and B.~Pfaff.
\emph{Revisiting the Open vSwitch Dataplane Ten Years Later}.
SIGCOMM, 2021.
\href{https://doi.org/10.1145/3452296.3472914}{DOI}.

\bibitem{ara2020}
G.~Ara, T.~Cucinotta, L.~Abeni and C.~Vitucci.
\emph{Comparative Evaluation of Kernel Bypass Mechanisms for High-performance Inter-container Communications}.
CLOSER, 2020.
\href{https://doi.org/10.5220/0009321200440055}{DOI}.

\end{thebibliography}

\end{document}
